\documentclass[10pt,a4paper]{amsart}
\usepackage[margin=19mm]{geometry}
\usepackage{amsmath,amssymb,mathtools}
\usepackage[scr=rsfs,cal=euler]{mathalfa}
\usepackage{xfrac}
\usepackage[unicode,hidelinks,bookmarks=false]{hyperref}
\newcommand{\ie}{i.e.\ }
\newcommand{\cf}{cf.\ }
\newcommand{\resp}{respectively\ }
\newcommand{\Subsection}{\subsection}
\providecommand{\nobreakdash}{\nobreak\hbox{-}}
\setlength{\emergencystretch}{2em}
\multlinegap0pt
\usepackage[shortlabels]{enumitem}
\usepackage{amssymb}
\usepackage{mathtools}
\usepackage{xcolor}
\newtheorem{corollary}{Corollary}
\newtheorem{lemma}{Lemma}
\newtheorem{theorem}{Theorem}
\usepackage{cleveref}
\crefname{corollary}{Corollary}{Corollarys}
\crefname{lemma}{Lemma}{Lemmas}
\crefname{theorem}{Theorem}{Theorems}
\newcommand{\N}{\mathbb N}
\newcommand{\R}{\mathbb R}
\newcommand{\Var}{\operatorname{Var}}
\newcommand{\abs}[1]{\left\lvert #1\right\rvert}
\newcommand{\cF}{\mathcal F}
\newcommand{\cG}{\mathcal G}
\newcommand{\cK}{\mathcal K}
\newcommand{\dd}{\,d}
\newcommand{\norm}[1]{\left\lVert #1\right\rVert}
\title{Quantitative Driver-Only Capacity Cores and Simultaneous Scalar It\^o Flows}
\author{Guangqian Zhao}
\date{Source: arXiv:2510.06054v6 (CC BY 4.0)}
\begin{document}
\maketitle

\noindent\emph{Source and licence.} Quantitative Driver-Only Capacity Cores and Simultaneous Scalar It\^o Flows. Version arXiv:2510.06054v6 (CC BY 4.0), distributed by arXiv under the Creative Commons Attribution 4.0 International licence, http://creativecommons.org/licenses/by/4.0/, as recorded in the arXiv licence metadata for this exact version and shown on the abstract page https://arxiv.org/abs/2510.06054v6. Full licence notice: \texttt{licenses/arxiv-2510.06054-CC-BY-4.0.txt}.\par\smallskip

\noindent\textit{Editorial scope.} Counted unit: the article's theorem thm:simultaneous-flow (source0.tex lines 2763--2869). The article's complete original proof of it is delivered with it (source0.tex lines 2871--3113, one \texttt{proof} environment of 243 lines). No mathematics is written, repaired or re-derived: the statement and the proof are the article's own lines, in the article's own notation, and no retained line is substituted or re-flowed.  Retained uncounted as the source-local context the proof needs: lines 1919--1922; lines 2279--2315; lines 2384--2386; lines 2392--2401; lines 2408--2445; lines 2412--2416; lines 2420--2425; lines 2429--2438; lines 2554--2573; lines 2607--2638; lines 2632--2637.  Nothing was omitted from any retained span, so the package carries no omission record.  No retained line contains a citation, so the wrapper carries no bibliography and no bibliography entry.  No retained line contains a figure environment, an image-inclusion command or drawing code, so no image is delivered and the package declares no asset dependency.  Editorial formatting only, in addition: an AMS \texttt{amsart} preamble that restates the article's own declarations and macros used by the retained lines, namely the environments \texttt{corollary}, \texttt{lemma}, \texttt{theorem} on separate counters, together with the standard packages those lines need.  The retained environments print in this document as Corollary 1, Lemma 1, Theorem 1 in order of appearance, whereas the article prints the same items with the numbering of its own full document.  The complete native source of the article as distributed in the arXiv e-print for this exact version is bundled unchanged as \texttt{sources/186578/source0.tex}.
\begin{equation}\label{eq:intrinsic-qv-identification}
  Q(X)=[X]^P
  \qquad P\text{-almost surely},
\end{equation}

\begin{corollary}[Parabolic diagonal]
\label{cor:parabolic-core}
Define
\begin{equation}\label{eq:parabolic-core}
  \cK_R:=\cK_{R,R^2},
  \qquad R\ge1.
\end{equation}
Then \((\cK_R)_{R\ge1}\) is increasing, and there are
\(c,C,C_*>0\), depending only on \((\eta,\gamma)\), such that
\begin{equation}\label{eq:core-capacity-rate}
  c_\Lambda(\cK_R^c)\le Ce^{-cR^2},
  \qquad R\ge1,
\end{equation}
\begin{equation}\label{eq:qv-core-rate}
  \sup_{x\in\cK_R}\norm{Q^n(x)-Q(x)}_\infty
  \le C_\gamma R^2\Lambda T2^{-\gamma n},
  \qquad n\ge1,
\end{equation}
\begin{equation}\label{eq:qv-uniform-on-core}
  \lim_{n\to\infty}
  \sup_{x\in\cK_R}\norm{Q^n(x)-Q(x)}_\infty=0,
\end{equation}
and
\begin{equation}\label{eq:qv-core-modulus}
  \norm{Q(x)-Q(\widetilde x)}_\infty
  \le C_{\eta,\gamma}R^{2-\alpha}
       (\Lambda T)^{1-\alpha/2}
       \norm{x-\widetilde x}_\infty^\alpha,
  \qquad x,\widetilde x\in\cK_R.
\end{equation}
The diagonal cores are stable under stopping and, for every \(s\in[0,T]\),
\begin{equation}\label{eq:core-concatenation}
  x,\widetilde x\in\cK_R
  \quad\Longrightarrow\quad
  x\otimes_s\widetilde x\in\cK_{C_*R}.
\end{equation}
\end{corollary}

\begin{equation}\label{eq:G-third-bound}
  \sup_{\theta\in\Theta}\norm{G_\theta'''}_\infty<\infty .
\end{equation}

\begin{align}
  Z_t^{\theta,s,y}(x,q)
  ={}&F_\theta(y)+x_t-x_s
  +\int_s^t\widehat b_\theta
       \bigl(Z_r^{\theta,s,y}(x,q)\bigr)\,\dd r
  \notag\\
  &+\int_s^t\widehat c_\theta
       \bigl(Z_r^{\theta,s,y}(x,q)\bigr)\,\dd q_r,
  \label{eq:Lamperti-equation}
\end{align}

\begin{lemma}[Deterministic parameterized flow]
\label{lem:deterministic-flow}
Equation \eqref{eq:Lamperti-equation} has a unique global continuous
solution.  On every set \(\{q_T\le A\}\), the map
\begin{equation}\label{eq:joint-flow-continuity}
  (x,q,\theta,s,t,y)
  \longmapsto
  \Phi_{s,t}^\theta(x,q;y)
\end{equation}
is jointly continuous, locally uniformly in \(y\), when \(x\) and \(q\)
carry the uniform topology.  It is causal and satisfies, for
\(s\le u\le t\),
\begin{equation}\label{eq:cocycle}
  \Phi_{s,t}^\theta(x,q;y)
  =
  \Phi_{u,t}^\theta
  \bigl(x,q;\Phi_{s,u}^\theta(x,q;y)\bigr).
\end{equation}
For fixed \((x,q,\theta,s,t)\), the map
\(y\mapsto\Phi_{s,t}^\theta(x,q;y)\) is a \(C^1\)-diffeomorphism of
\(\R\), and it preserves orientation.  More precisely,
\begin{equation}\label{eq:flow-derivative}
  \partial_y\Phi_{s,t}^\theta(x,q;y)
  =
  \frac{\sigma_\theta(\Phi_{s,t}^\theta(x,q;y))}
       {\sigma_\theta(y)}
  \exp\!\left\{
    \int_s^t\widehat b_\theta'(Z_r)\,\dd r
    +\int_s^t\widehat c_\theta'(Z_r)\,\dd q_r
  \right\}>0.
\end{equation}
Consequently, on \(\{q_T\le A\}\) there are constants
\(0<c_A\le C_A<\infty\), independent of
\((x,\theta,s,t,y)\), such that
\begin{equation}\label{eq:flow-bi-Lipschitz}
  c_A\le\partial_y\Phi_{s,t}^\theta(x,q;y)\le C_A.
\end{equation}
\end{lemma}

\begin{equation}
  (x,q,\theta,s,t,y)
  \longmapsto
  \Phi_{s,t}^\theta(x,q;y)
\end{equation}

\begin{equation}
  \Phi_{s,t}^\theta(x,q;y)
  =
  \Phi_{u,t}^\theta
  \bigl(x,q;\Phi_{s,u}^\theta(x,q;y)\bigr).
\end{equation}

\begin{equation}
  \partial_y\Phi_{s,t}^\theta(x,q;y)
  =
  \frac{\sigma_\theta(\Phi_{s,t}^\theta(x,q;y))}
       {\sigma_\theta(y)}
  \exp\!\left\{
    \int_s^t\widehat b_\theta'(Z_r)\,\dd r
    +\int_s^t\widehat c_\theta'(Z_r)\,\dd q_r
  \right\}>0.
\end{equation}

\begin{lemma}[Quantitative Stieltjes stability]
\label{lem:quantitative-stieltjes}
Fix \(0<\eta<1\).  Let \(q,\widetilde q\in\mathcal A_\Lambda\), and let
\(f\in C([0,T])\) satisfy
\[
  \norm{f}_\infty\le B,
  \qquad
  \abs{f(t)-f(s)}\le H\abs{t-s}^\eta.
\]
If \(\delta=\norm{q-\widetilde q}_\infty\), then
\begin{equation}\label{eq:quantitative-stieltjes}
  \sup_{0\le s\le t\le T}
  \left|
    \int_s^t f(r)\,\dd q_r
    -\int_s^t f(r)\,\dd\widetilde q_r
  \right|
  \le C_{B,H,\eta,\Lambda,T}
  \bigl(\delta+\delta^\eta\bigr).
\end{equation}
\end{lemma}

\begin{lemma}[Uniform weighted quadratic variation on the cores]
\label{lem:uniform-weighted-qv}
Fix \(R,m<\infty\).  Uniformly over
\[
  x\in\cK_R,\quad \theta\in\Theta,\quad
  s\in[0,T],\quad \abs{y}\le m,
\]
the quadratic-variation sums of
\(Z^{\theta,s,y}(x,Q(x))\) along \(\pi_n^s\) define positive discrete
measures on \([s,T]\) which converge weakly to the restriction of
\(\dd Q(x)\) to \([s,T]\), uniformly over the indicated parameters.  More
precisely, if \(\mathscr F\) is any uniformly bounded and equicontinuous
family of continuous weights, then
\begin{equation}\label{eq:uniform-weighted-qv}
  \sup
  \left|
    \sum_{\substack{[u,v]\in\pi_n^s\\v\le t}}
      f(u)\bigl(Z_v-Z_u\bigr)^2
    -
    \int_s^t f(r)\,\dd Q_r(x)
  \right|
  \longrightarrow0,
\end{equation}
where the supremum is also over \(t\in[s,T]\) and \(f\in\mathscr F\).
Furthermore,
\begin{equation}\label{eq:uniform-cubic-remainder}
  \sup
  \sum_{[u,v]\in\pi_n^s}
  \abs{Z_v-Z_u}^3
  \longrightarrow0.
\end{equation}
\end{lemma}

\begin{equation}
  \sup
  \sum_{[u,v]\in\pi_n^s}
  \abs{Z_v-Z_u}^3
  \longrightarrow0.
\end{equation}

\begin{theorem}[Simultaneous driver-only It\^o flow]
\label{thm:simultaneous-flow}
For \(x\in\cG_{\mathrm{dyad}}\), put
\begin{equation}\label{eq:driver-flow}
  S_{s,t}^\theta(x;y)
  :=\Phi_{s,t}^\theta(x,Q(x);y).
\end{equation}
When \(S(X)\) is viewed as a random field on all of \(\Omega\), we extend it
by the identity flow \(S_{s,t}^\theta(x;y)=y\) on
\(\cG_{\mathrm{dyad}}^c\).  This yields a total Borel random field.
Then the following statements hold.

\begin{enumerate}[label=(\roman*),leftmargin=2.2em]
  \item For every \(x\in\cG_{\mathrm{dyad}}\), every
  \((\theta,s,y)\in\Theta\times[0,T]\times\R\), and every \(t\ge s\),
  the left sums
  \begin{equation}\label{eq:Follmer-sums}
    J_{s,t}^{n,\theta}(x;y)
    =
    \sum_{[u,v]\in\pi_n^s\,:\,s\le u<t}
    \sigma_\theta\bigl(S_{s,u}^\theta(x;y)\bigr)
    \bigl(x_{v\wedge t}-x_u\bigr)
  \end{equation}
  converge uniformly in \(t\).  Calling the limit
  \(J_{s,t}^\theta(x;y)\), one has the deterministic It\^o equation
  \begin{equation}\label{eq:pathwise-sde-flow}
    S_{s,t}^\theta(x;y)
    =y+
    \int_s^t b_\theta\bigl(S_{s,r}^\theta(x;y)\bigr)\,\dd r
    +J_{s,t}^\theta(x;y).
  \end{equation}
  The convergence and the identity hold for all parameters on the single
  domain \(\cG_{\mathrm{dyad}}\).

  \item The field
  \begin{equation}\label{eq:flow-valued-map}
    x\longmapsto
    \left(
      (\theta,s,t,y)\longmapsto S_{s,t}^\theta(x;y)
    \right)
  \end{equation}
  and the corresponding integral field
  \[
    x\longmapsto
    \left(
      (\theta,s,t,y)\longmapsto J_{s,t}^\theta(x;y)
    \right)
  \]
  are Borel and causal from \(\cG_{\mathrm{dyad}}\) into
  \[
    C_{\mathrm{loc}}
    \bigl(\Theta\times\{(s,t):0\le s\le t\le T\}\times\R;\R\bigr).
  \]
  Here \(C_{\mathrm{loc}}\) carries the compact-open topology.
  The solution field satisfies the cocycle identity \eqref{eq:cocycle} and is a
  \(C^1\)-diffeomorphic flow in the initial value.  The integral field
  satisfies, for \(s\le u\le t\),
  \begin{equation}\label{eq:integral-cocycle}
    J_{s,t}^\theta(x;y)
    =J_{s,u}^\theta(x;y)
     +J_{u,t}^\theta
       \bigl(x;S_{s,u}^\theta(x;y)\bigr).
  \end{equation}

  \item For the compact sets \(\cK_R\) in
  \Cref{cor:parabolic-core}, the flow-valued map
  \eqref{eq:flow-valued-map} and the integral-valued map are uniformly
  continuous on \(\cK_R\).  In fact, with
  \begin{equation}\label{eq:flow-holder-exponent}
    \beta_{\mathrm{flow}}
    =\eta\alpha
    =\frac{\eta\gamma}{1-\eta+\gamma},
  \end{equation}
  there is \(C_R<\infty\) such that, whenever
  \(x,\widetilde x\in\cK_R\) and
  \(e=\norm{x-\widetilde x}_\infty\le1\),
  \begin{align}
    \sup_{\substack{\theta\in\Theta,\ 0\le s\le t\le T\\y\in\R}}
    \Bigl(&
      \abs{S_{s,t}^\theta(x;y)-S_{s,t}^\theta(\widetilde x;y)}
      \notag\\[-0.2em]
      &+
      \abs{J_{s,t}^\theta(x;y)-J_{s,t}^\theta(\widetilde x;y)}
    \Bigr)
    \le C_Re^{\beta_{\mathrm{flow}}}.
    \label{eq:flow-raw-holder}
  \end{align}
  In addition, the sums \eqref{eq:Follmer-sums} converge locally uniformly in
  \((x,\theta,s,t,y)\) on
  \(\cK_R\times\Theta\times\{s\le t\}\times\R\).

  \item Under every \(P\in\mathfrak M_\Lambda\), for every fixed
  \((\theta,s,y)\), the process
  \(S_{s,\cdot}^\theta(X;y)\) is the unique strong solution of
  \begin{equation}\label{eq:classical-scalar-SDE}
    Y_t
    =y+\int_s^t b_\theta(Y_r)\,\dd r
       +\int_s^t\sigma_\theta(Y_r)\,\dd X_r,
    \qquad t\ge s.
  \end{equation}
  For every fixed \((\theta,s,y)\), the corresponding section of \(J\) is a
  version of the It\^o integral in \eqref{eq:classical-scalar-SDE}.  Its
  pathwise equation and cocycle identities hold on
  \(\cG_{\mathrm{dyad}}\) for all parameters, while the stochastic
  identifications hold up to indistinguishability for each fixed parameter.
\end{enumerate}
\end{theorem}

\begin{proof}
Fix \(x\in\cG_{\mathrm{dyad}}\), write \(q=Q(x)\), and abbreviate
\(Z=Z^{\theta,s,y}(x,q)\), \(Y=G_\theta(Z)\).  From
\eqref{eq:Lamperti-equation}, put
\[
  Z_t=Z_s+x_t-x_s+A_t,
\]
where \(A_s=0\) and \(A\) is continuous with finite variation.  Let
\(\Delta_n=T2^{-n}\), and, for a continuous path \(w\), define
\[
  V_{n,s}^w(t)
  =
  \sum_{\substack{[u,v]\in\pi_n^s\\s\le u<t}}
    \bigl(w_{v\wedge t}-w_u\bigr)^2,
  \qquad s\le t\le T.
\]
Comparing the partition obtained by inserting \(s\) with the original
dyadic partition gives
\begin{equation}\label{eq:all-path-inserted-qv}
  \sup_{s\le t\le T}
  \left|
    V_{n,s}^x(t)-(q_t-q_s)
  \right|
  \le
  2\norm{Q^n(x)-q}_\infty+4\omega_x(\Delta_n)^2
  \longrightarrow0.
\end{equation}
Here \(\omega_x\) is the modulus of continuity of \(x\).  Thus
\eqref{eq:all-path-inserted-qv} holds for every
\(x\in\cG_{\mathrm{dyad}}\), independently of the compact cores.

The finite-variation terms satisfy, uniformly in \(t\in[s,T]\),
\[
  \left|V_{n,s}^Z(t)-V_{n,s}^x(t)\right|
  \le
  2\omega_x(\Delta_n)\Var(A;[s,T])
  +\omega_A(\Delta_n)\Var(A;[s,T])
  \longrightarrow0.
\]
Consequently \(V_{n,s}^Z\to q-q_s\) uniformly.  Moreover,
\[
  \sup_{s\le t\le T}
  \left|
    \mu_{n,s}^Z([s,t])-V_{n,s}^Z(t)
  \right|
  \le \omega_Z(\Delta_n)^2\longrightarrow0.
\]
Therefore the positive measures
\[
  \mu_{n,s}^Z
  =\sum_{[u,v]\in\pi_n^s}(Z_v-Z_u)^2\delta_v
\]
have uniformly bounded mass and cumulative functions converging uniformly to
\(q-q_s\).  Step-function approximation of the continuous weight
\(G_\theta''(Z)\), together with the vanishing unfinished-cell error, yields
\begin{equation}\label{eq:all-path-weighted-qv}
  \sup_{s\le t\le T}
  \left|
    \sum_{\substack{[u,v]\in\pi_n^s\\u<t}}
      G_\theta''(Z_u)(Z_{v\wedge t}-Z_u)^2
    -\int_s^tG_\theta''(Z_r)\,\dd q_r
  \right|
  \longrightarrow0.
\end{equation}

Taylor's formula on the intervals of \(\pi_n^s\) gives a remainder
\(R_n(t)\) satisfying, by \eqref{eq:G-third-bound},
\[
  \sup_{s\le t\le T}\abs{R_n(t)}
  \le
  C\omega_Z(\Delta_n)
   \sup_{s\le t\le T}V_{n,s}^Z(t)
  \longrightarrow0.
\]
Together with \eqref{eq:all-path-weighted-qv}, this proves uniform convergence
of the F\"ollmer left sums of \(G_\theta'(Z)\) against \(Z\).  The ordinary
left sums against \(A\) also converge uniformly, since \(A\) has finite
variation.  Subtracting them leaves precisely the sums
\eqref{eq:Follmer-sums}.  Their limit is
\begin{align*}
  Y_t-y
  &-\int_s^tG_\theta'(Z_r)\widehat b_\theta(Z_r)\,\dd r\\
  &-\int_s^t
    \left(
      G_\theta'(Z_r)\widehat c_\theta(Z_r)
      +\frac12G_\theta''(Z_r)
    \right)\dd q_r.
\end{align*}
Since
\[
  G_\theta'(Z_r)\widehat b_\theta(Z_r)
    =b_\theta(Y_r),
  \qquad
  G_\theta'(Z_r)\widehat c_\theta(Z_r)
    +\frac12G_\theta''(Z_r)=0,
\]
the limit equals \(Y_t-y-\int_s^t b_\theta(Y_r)\,\dd r\).  This proves the
uniform-in-\(t\) convergence in (i) and \eqref{eq:pathwise-sde-flow} for all
parameters on \(\cG_{\mathrm{dyad}}\).

Set
\[
  \Delta_T=\{(s,t):0\le s\le t\le T\},
  \qquad
  \mathsf P_0=\Theta\times\Delta_T\times\R,
\]
and exhaust \(\mathsf P_0\) by the compact sets
\(\mathsf P_m=\Theta\times\Delta_T\times[-m,m]\).  On every set
\(\{q_T\le a\}\), \Cref{lem:deterministic-flow} and compactness of
\(\mathsf P_m\) show that
\[
  (x,q)\longmapsto
  \left(
    (\theta,s,t,y)\longmapsto\Phi_{s,t}^\theta(x,q;y)
  \right)
\]
is continuous into \(C_{\mathrm{loc}}(\mathsf P_0;\R)\).  Since
\[
  C_\uparrow([0,T])
  =\bigcup_{a\in\N}\{q\in C_\uparrow([0,T]):q_T\le a\}
\]
is a countable union of closed sets, the same function-space-valued map is
Borel on \(\Omega\times C_\uparrow([0,T])\).  Composing it with the Borel
map \(x\mapsto(x,Q(x))\) proves that the solution field in
\eqref{eq:flow-valued-map} is a
\(C_{\mathrm{loc}}(\mathsf P_0;\R)\)-valued Borel map.  Causality follows
from the causality of \(Q\) and of the deterministic equation.
Because \(\cG_{\mathrm{dyad}}\) is Borel and the identity flow is a fixed
element of the target space, the extension on \(\cG_{\mathrm{dyad}}^c\) is
Borel as well.

Equation \eqref{eq:pathwise-sde-flow} gives the identity
\begin{equation}\label{eq:J-from-S}
  J_{s,t}^\theta(x;y)
  =S_{s,t}^\theta(x;y)-y
   -\int_s^t b_\theta(S_{s,r}^\theta(x;y))\,\dd r.
\end{equation}
The operator on \(C_{\mathrm{loc}}(\mathsf P_0;\R)\) defined by the
right-hand side of \eqref{eq:J-from-S} is continuous; on each
\(\mathsf P_m\), its Lipschitz constant is at most \(1+LT\).  Hence the
integral field is also Borel and causal.  Splitting
\eqref{eq:J-from-S} at \(u\) and using the solution cocycle proves
\eqref{eq:integral-cocycle}.  Formula
\eqref{eq:flow-derivative} gives the \(C^1\)-diffeomorphism assertion and
completes (ii).

On \(\cK_R\), the map
\(x\mapsto(x,Q(x))\) is continuous and has compact image.  Composition with
\eqref{eq:joint-flow-continuity} is therefore continuous into the stated
local-uniform function space and hence uniformly continuous; the same is
true for \(J\) by \eqref{eq:J-from-S}.

We prove the stronger modulus \eqref{eq:flow-raw-holder}.  Take
\(x,\widetilde x\in\cK_R\), put
\(q=Q(x)\), \(\widetilde q=Q(\widetilde x)\), and write
\(\delta=\norm{q-\widetilde q}_\infty\).  For fixed
\((\theta,s,y)\), let \(Z,\widetilde Z\) denote the two Lamperti solutions.
Write again
\[
  B_b=\sup_{\theta\in\Theta}\norm{\widehat b_\theta}_\infty,
  \qquad
  B_c=\sup_{\theta\in\Theta}\norm{\widehat c_\theta}_\infty.
\]
The boundedness of \(\widehat b,\widehat c\), the \(\eta\)-H\"older bound on
\(\widetilde x\), and \(\widetilde q\in\mathcal A_\Lambda\) give, uniformly
in all parameters,
\[
  \abs{\widetilde Z_t-\widetilde Z_u}
  \le R\sqrt\Lambda\,T^{1/2-\eta}\abs{t-u}^\eta
      +(B_b+B_c\Lambda)\abs{t-u}
  \le C_R\abs{t-u}^\eta.
\]
Thus \(r\mapsto\widehat c_\theta(\widetilde Z_r)\) belongs to a uniformly
bounded, uniformly \(\eta\)-H\"older family.  Subtracting the two Lamperti
equations, applying \Cref{lem:quantitative-stieltjes} to the term driven by
\(q-\widetilde q\), and then applying Gronwall's inequality for the measure
\(\dd r+\dd q_r\), yields
\[
  \sup_{t\in[s,T]}\abs{Z_t-\widetilde Z_t}
  \le C_R\bigl(e+\delta+\delta^\eta\bigr),
  \qquad e=\norm{x-\widetilde x}_\infty.
\]
Since the maps \(G_\theta\) are uniformly Lipschitz and
\(\delta\le C_Re^\alpha\) by \eqref{eq:qv-core-modulus}, we obtain
\[
  \sup_{\substack{\theta\in\Theta,\ 0\le s\le t\le T\\y\in\R}}
  \abs{S_{s,t}^\theta(x;y)-S_{s,t}^\theta(\widetilde x;y)}
  \le C_Re^{\eta\alpha},
  \qquad e\le1.
\]
The same estimate for \(J\) follows from \eqref{eq:J-from-S} and the common
Lipschitz bound on \(b_\theta\).  This proves \eqref{eq:flow-raw-holder}.

The local uniform assertion for the sums follows
from the quantitative deterministic proof of F\"ollmer's formula.  On
\(\cK_R\times\Theta\times\{s\le t\}\times[-m,m]\), the families \(x\),
\(q=Q(x)\), \(Z\), and their finite-variation parts are uniformly bounded
and equicontinuous; the latter also have uniformly bounded variation.
Taylor's formula separates four errors.  First, the finite-variation left
sums converge uniformly because
\[
  \sup_t\left|
    \sum_{\substack{[u,v]\in\pi_n^s\\u<t}}
      G_\theta'(Z_u)(A_{v\wedge t}-A_u)
    -\int_s^tG_\theta'(Z_r)\,\dd A_r
  \right|
  \le
  \omega_{G_\theta'(Z)}(T2^{-n})\Var(A;[s,T])
  \longrightarrow0.
\]
Second, the uniform third-derivative bound
\eqref{eq:G-third-bound} and
\Cref{lem:uniform-weighted-qv}, applied to the equicontinuous family of
weights \(G_\theta''(Z)\), give uniform convergence of the quadratic terms.
Third, the cubic Taylor remainders vanish by
\eqref{eq:uniform-cubic-remainder}.  Fourth, the sole unfinished terminal
cell is controlled by the common oscillation modulus of \(Z\).  These four
estimates prove the local uniform convergence in (iii).

Finally fix \(P\in\mathfrak M_\Lambda\) and an adapted continuous version of
\([X]^P\).  Solve \eqref{eq:Lamperti-equation} with the two adapted inputs
\((X,[X]^P)\).  Causality of the deterministic equation, or equivalently its
Picard iteration, shows that, for every fixed \((\theta,s,y)\), the resulting
processes \(\overline Z\) and \(\overline S=G_\theta(\overline Z)\) are
adapted.  By
\eqref{eq:intrinsic-qv-identification},
\[
  \overline S_{s,\cdot}^\theta(y)
  =S_{s,\cdot}^\theta(X;y)
  \qquad P\text{-almost surely}.
\]
The equality holds simultaneously in the parameters on the single
\(P\)-full event \(\{Q(X)=[X]^P\}\).  Under the usual augmentation, its
complement belongs to \(\cF_0^P\).  Hence, for every fixed \((\theta,s,y)\),
the section \(S_{s,\cdot}^\theta(X;y)\) is adapted under \(P\); continuity
then implies progressive measurability.  The process
\(\overline Z\) is a continuous semimartingale.  Applying the classical
It\^o formula to \(G_\theta(\overline Z)\) and using the two cancellation
identities proves \eqref{eq:classical-scalar-SDE}.  Global Lipschitz
continuity gives strong uniqueness.  Equation \eqref{eq:J-from-S} then
identifies \(J\), up to indistinguishability, with the It\^o integral for each
fixed \((\theta,s,y)\), proving (iv).
\end{proof}

\end{document}
